% ============================================================
% ANEXO A: Códigos Fuente
% ============================================================

\Anexo{Códigos Fuente del Sistema}

Este anexo contiene los códigos fuente desarrollados para la implementación del gemelo digital. El código está organizado en tres módulos principales que corresponden a los componentes de hardware y software del sistema.

\subsection*{A.1 Código ESP32 - Nodo Ejecutor}

El siguiente código corresponde al firmware del microcontrolador ESP32 encargado de controlar los actuadores (electroválvulas, pistones y servomotor). El código implementa la lógica de recepción de comandos MQTT y la activación de los relés correspondientes.

\begin{verbatim}
Archivo: Anexos/Codigos/esp32_ejecutor.py
\end{verbatim}

\subsection*{A.2 Código ESP32 - Nodo Observador}

El siguiente código corresponde al firmware del microcontrolador ESP32 encargado de la lectura de sensores (infrarrojos, ultrasónico y transductor de presión). El código implementa el filtrado de señales y la publicación periódica de telemetría vía MQTT.

\begin{verbatim}
Archivo: Anexos/Codigos/esp32_observador.py
\end{verbatim}

\subsection*{A.3 Código FSM - Webots (Python)}

El siguiente código implementa la Máquina de Estados Finitos (FSM) en Python dentro del entorno de simulación Webots. Este script controla la cinemática del gemelo digital y se sincroniza con el hardware físico mediante MQTT.

\begin{verbatim}
Archivo: Anexos/Codigos/fsm_webots.py
\end{verbatim}

% ============================================================
% ANEXO B: Diagramas Eléctricos
% ============================================================

\Anexo{Diagramas Eléctricos y Esquemáticos}

Este anexo contiene los diagramas eléctricos desarrollados en KiCad para la implementación de la caja de control externa. Los esquemáticos incluyen el diseño de la etapa de aislamiento galvánico, la conexión de sensores y la distribución de pines del ESP32.

\begin{verbatim}
Archivos:
- Anexos/Diagramas/esquematico_kiCad.pdf
- Anexos/Diagramas/pcb_kiCad.pdf
\end{verbatim}

% ============================================================
% ANEXO C: Modelos 3D
% ============================================================

\Anexo{Modelos 3D Diseñados en FreeCAD}

Este anexo presenta los modelos tridimensionales diseñados en FreeCAD para la construcción del gemelo digital en Webots. Los modelos respetan las dimensiones y restricciones mecánicas de la máquina real.

\subsection*{C.1 Modelo de la Tolva de Alimentación}

\begin{verbatim}
Archivo: Anexos/Modelos_3D/tolva.fcstd
\end{verbatim}

\subsection*{C.2 Modelo de la Trampilla de Dosificación}

\begin{verbatim}
Archivo: Anexos/Modelos_3D/trampilla.fcstd
\end{verbatim}

\subsection*{C.3 Modelo del Mecanismo de Sellado}

\begin{verbatim}
Archivo: Anexos/Modelos_3D/sellador.fcstd
\end{verbatim}

\subsection*{C.4 Modelo de la Estructura General}

\begin{verbatim}
Archivo: Anexos/Modelos_3D/estructura.fcstd
\end{verbatim}

\subsection*{C.5 Proyecto Completo en Webots}

El modelo cinemático completo fue integrado en Webots, donde se programaron las restricciones de movimiento y las interacciones físicas.

\begin{verbatim}
Archivo: Anexos/Webots/proyecto_gemelo.wbt
\end{verbatim}

% ============================================================
% ANEXO D: Configuración SCADA y Node-RED
% ============================================================

\Anexo{Configuración del SCADA en Node-RED}

Este anexo documenta la configuración del panel SCADA desarrollado en Node-RED. Se incluye el flujo completo de nodos, la configuración de tópicos MQTT y la estructura del dashboard.

\begin{verbatim}
Archivo: Anexos/Node-RED/flujo_node-red.json
\end{verbatim}

\subsection*{D.1 Configuración de Tópicos MQTT}

\begin{verbatim}
maquina/sensor/presion
maquina/sensor/nivel_tolva
maquina/sensor/distancia
maquina/actuador/electrovalvula
maquina/actuador/piston
maquina/simulacion/actuador/#
maquina/estado/fsm
\end{verbatim}

% ============================================================
% ANEXO E: Resultados de Pruebas
% ============================================================

\Anexo{Resultados de las Pruebas de Validación}

Este anexo contiene los resultados detallados de las cinco pruebas de validación realizadas. Los datos incluyen métricas de OEE, disponibilidad, rendimiento y latencia MQTT.

\subsection*{E.1 Reportes de Pruebas}

\begin{verbatim}
Archivos:
- Anexos/Resultados/reporte_prueba1.pdf
- Anexos/Resultados/reporte_prueba2.pdf
- Anexos/Resultados/reporte_prueba3.pdf
- Anexos/Resultados/reporte_prueba4.pdf
- Anexos/Resultados/reporte_prueba5.pdf
\end{verbatim}

\subsection*{E.2 Base de Datos}

El archivo SQLite contiene el registro histórico de todas las variables de proceso, transiciones de estado y eventos del sistema.

\begin{verbatim}
Archivo: Anexos/Datos/datos_pruebas.sqlite
\end{verbatim}
